\documentclass[12pt]{article}
\usepackage[utf8]{inputenc}
\usepackage[T1]{fontenc}
\usepackage{hyperref}
\title{Woria Language Resources for K-12 Curriculum Development}
\author{Compiled by Claude Sonnet 5.5 (claude-sonnet-5-5) for Andr\'as Kornai}
\date{}
\begin{document}
\maketitle

\section{Introduction}
Woria (ISO 639-3: wor; Glottocode: wori1246) is a nearly extinct language of Botawa village, Waropen Regency, on the eastern shore of Cenderawasih (Geelvink) Bay, Papua, Indonesia. Estimates give about 5 or 6 speakers. UNESCO classes it as critically endangered. The village also hosts the larger Demisa and Waropen languages, which put pressure on Woria, and Woria has no literacy tradition. It is lexically similar to the East Geelvink Bay languages and presumably belongs to that family, but it is too poorly attested to be sure.

\section{Primary Educational Resources}
None found. No primers, textbooks or school programmes exist in the sources I saw.

\section{Linguistic and Reference Materials}
\begin{itemize}
\item Wikipedia article on the Woria language: \url{https://en.wikipedia.org/wiki/Woria_language}
\item Endangered Languages Project entry: \url{https://endangeredlanguages.com/lang/2803}
\item Verbix endangered-languages page for Woria: \url{https://docs.verbix.com/EndangeredLanguages/Woria}
\end{itemize}

\section{Religious and Community Resources}
None found.

\section{Digital and Audio Resources}
According to Wikipedia, Paradisec holds an open-access collection from the linguist Bert Voorhoeve that includes Woria material. I did not see a direct URL for it in search results, so it is not in the JSON list.

\section{Curriculum Development Context}
With so few speakers and almost no documentation, any work would be language documentation and revitalisation, not a K-12 curriculum. The first steps would be recording the remaining speakers and locating the Paradisec material.
\end{document}